
\section{$\OO_X$ 模的导出函子上同调}

我们想把导出函子的这套机器用起来，为概形 $X$ 上的\trQCSheaf{}定义上同调。与其在范畴 $\mathit{QCoh}_X$ 中工作，出于若干理由，在更大的范畴 $\mathit{Mod}_{\OO_X}$ 中工作会更简单（但参看\mbox{不重要注记~23.5.7}）。

\textbf{23.4.1. \trhTheorem。}\ ——\ \emph{设 $(X,\OO_X)$ 是一个\trRingedSpace。则 $\OO_X$-\trModule{}的范畴 $\mathit{Mod}_{\OO_X}$ 有足够多内射对象。}

作为一项副产品（对别人比对我们更有用），取 $\OO_X=\mathbb{Z}$，我们看到：固定拓扑空间上阿贝尔群层的范畴有足够多内射对象。

\textbf{23.4.2.}\ \emph{证明。}\quad 我们沿循 Godement 的做法 [GrS, pp.~27--28]，用一系列练习来证明\mbox{定理~23.4.1}。设 $\mathscr{F}$ 是一个 $\OO_X$-\trModule。我们将展示一个单射 $\mathscr{F}\hookrightarrow\mathscr{Q}'$，其中 $\mathscr{Q}'$ 是某个内射 $\OO_X$-\trModule。对每个 $\mathfrak{p}\in X$，选取一个包含 $\mathscr{F}_\mathfrak{p}\hookrightarrow Q_\mathfrak{p}$，其中 $Q_\mathfrak{p}$ 是一个内射 $\OO_{X,\mathfrak{p}}$-\trModule（这是可能的，因为 $\OO_{X,\mathfrak{p}}$-\trModule{}的范畴有足够多内射对象，\mbox{练习~23.2.K}）。

\textbf{23.4.A.（练习）}\quad 证明：\trSkyscraperSheaf{} $\mathscr{Q}_\mathfrak{p}:=i_{\mathfrak{p},\ast}Q_\mathfrak{p}$（带有位于点 $\mathfrak{p}\in X$ 处的\trModule{} $Q_\mathfrak{p}$）是一个内射 $\OO_X$-\trModule。（你可以用\trAbstractNonsense{}漂亮地做到这一点——动用\mbox{练习~23.5.B}——但动手去做也同样快。）

\textbf{23.4.B.（容易练习）}\quad 证明：\trAbelianCategory{}中内射对象的（可能是无限的）乘积仍是内射的。

由前两道练习，$\mathscr{Q}':=\prod_{\mathfrak{p}\in X}\mathscr{Q}_\mathfrak{p}$ 是一个内射 $\OO_X$-\trModule。

\textbf{23.4.C.（容易练习）}\quad 通过审视茎，证明自然映射 $\mathscr{F}\to\mathscr{Q}'$ 是一个单射。

这样，\mbox{定理~23.4.1} 的证明便告完成。\hfill$\square$

我们现在可以给出好几个定义了。

\textbf{23.4.3.}\ \emph{\trhDefinition。}\quad 若 $(X,\OO_X)$ 是一个\trRingedSpace，而 $\mathscr{F}$ 是一个 $\OO_X$-\trModule，则定义 $\mathrm{H}^i(X,\mathscr{F})$ 为 $\mathrm{R}^i\Gamma(X,\mathscr{F})$。此外，若 $\pi\colon(X,\OO_X)\to(Y,\OO_Y)$ 是\trRingedSpace{}之间的一个映射，则我们有\textbf{\trHigherPushforward}（或\textbf{导出前推}）$\mathrm{R}^i\pi_\ast\colon\mathit{Mod}_{\OO_X}\to\mathit{Mod}_{\OO_Y}$。

这些概念我们早先在特殊情形中定义过：第 18 章中，对拟紧分离概形上的\trQCSheaf{}（就 $\mathrm{H}^i$ 而言），或对概形的拟紧分离态射（就 $\mathrm{R}^i\pi_\ast$ 而言）。我们很快将在 \S23.5 中证明，这些较早的定义与\mbox{定义~23.4.3} 相一致。因此，导出函子的定义比我们的切赫定义适用面广泛得多。但值得指出的是，我们将要用到的几乎所有东西都可以从切赫定义中得出。一个值得注意的例外是\trLeray{}\trSpectralSequence，我们现在就来讨论它。

\textbf{23.4.4. \trLeray{}谱序列。}\quad

\textbf{23.4.5. \trhTheorem（\trLeray{}谱序列）。}\ ——\ \emph{设 $\pi\colon(X,\OO_X)\to(Y,\OO_Y)$ 是\trRingedSpace{}之间的一个\trMorphism。对任何 $\OO_X$-\trModule{} $\mathscr{F}$，存在一个谱序列，其 $E_2$ 项由 $\mathrm{H}^q(Y,\mathrm{R}^p\pi_\ast\mathscr{F})$ 给出，抵达 $\mathrm{H}^{p+q}(X,\mathscr{F})$。}

这是\trGrothendieck 函子复合谱序列（\mbox{定理~23.3.5}）的一个直接推论——只须证明内射 $\OO$-\trModule{}的前推是\trAcyclic{}的 $\OO$-\trModule。我们现在就来证明这件事。

\textbf{23.4.6.}\ \emph{\trhDefinition。}\quad 我们给出一个本身也很重要的中间定义。拓扑空间上的层 $\mathscr{F}$ 称作\textbf{\trFlasque}的（有时也称作 flabby），如果一切限制映射都是满射，也就是说，如果 $\mathrm{res}_{\mathfrak{U}\subset\mathfrak{V}}\colon\mathscr{F}(\mathfrak{V})\to\mathscr{F}(\mathfrak{U})$ 对一切 $\mathfrak{U}\subset\mathfrak{V}$ 都是满射。

\textbf{23.4.D.（练习）}\quad 设 $(X,\OO_X)$ 是一个\trRingedSpace。
（a）证明：若
\[
\tag{23.4.6.1}\label{eq:flasque-ses}
0 \longrightarrow \mathscr{F}' \longrightarrow \mathscr{F} \longrightarrow \mathscr{F}'' \longrightarrow 0
\]
是 $\OO_X$-\trModule{}的一个正合序列，而 $\mathscr{F}'$ 是\trFlasque{}的，则 (23.4.6.1) 在任意开集 $\mathfrak{U}$ 的截面上正合，即 $0\to\mathscr{F}'(\mathfrak{U})\to\mathscr{F}(\mathfrak{U})\to\mathscr{F}''(\mathfrak{U})\to 0$
是正合的。
（b）给定正合序列 \eqref{eq:flasque-ses}，若 $\mathscr{F}'$ 是\trFlasque{}的，证明：$\mathscr{F}$ 是\trFlasque{}的当且仅当 $\mathscr{F}''$ 是\trFlasque{}的。

\textbf{23.4.E.（容易练习：松软层的前推仍是松软层）}\quad
（a）设 $\pi\colon X\to Y$ 是拓扑空间之间的一个连续映射，而 $\mathscr{F}$ 是 $X$ 上的一个由集合构成的\trFlasque{}层。证明 $\pi_\ast\mathscr{F}$ 是 $Y$ 上的一个\trFlasque{}层。
（b）设 $\pi\colon(X,\OO_X)\to(Y,\OO_Y)$ 是\trRingedSpace{}之间的一个\trMorphism，而 $\mathscr{F}$ 是一个\trFlasque{}的 $\OO_X$-\trModule。证明 $\pi_\ast\mathscr{F}$ 是一个\trFlasque{}的 $\OO_Y$-\trModule。

\textbf{23.4.7.}\ \emph{零扩张：逆像函子的一个偶尔扮演的左伴随。}\quad 为了证明内射层是\trFlasque{}的（\mbox{练习~23.4.H}），我们引入「\textbf{\trExtensionByZero}函子」。除了总是一个左伴随之外，当 $\pi$ 是开子集的包含时，$\pi^{-1}$ 有时还可以是一个右伴随。设 $i\colon\mathfrak{U}\hookrightarrow Y$ 是开子集到 $Y$ 中的一个包含。如下定义\textbf{\trExtensionByZero}函子 $i_!\colon\mathit{Mod}_{\OO_\mathfrak{U}}\to\mathit{Mod}_{\OO_Y}$。设 $\mathscr{F}$ 是一个 $\OO_\mathfrak{U}$-\trModule。对 $W\subset Y$，定义 $(i^{\mathrm{pre}}\mathscr{F})(W)=\mathscr{F}(W)$（若 $W\subset\mathfrak{U}$），否则为 $0$（限制映射取显然的那些）。这显然是一个预层 $\OO_Y$-\trModule。定义 $i_!$ 为 $(i^{\mathrm{pre}})^{\mathrm{sh}}$。注意 $i_!\mathscr{F}$ 是一个 $\OO_Y$-\trModule，并且这定义了一个函子。（符号「$!$」读作「shriek」（尖叫），大概是因为人们看到它就会尖叫。于是 $i_!$ 读作「i-lower-shriek」。）

\textbf{23.4.F.（练习）}\quad 证明：$i^{\mathrm{pre}}\mathscr{F}$ 满足恒等公理，但未必是一个层。（我们后面不会用到这一点，但它或许能让你看出这为什么叫「零扩张」。可能的例子来源：$\mathbb{R}$ 上的连续函数。）

\textbf{23.4.G.（练习）}\quad
（a）对 $\mathfrak{q}\in Y$，证明 $(i_!\mathscr{F})_\mathfrak{q}=\mathscr{F}_\mathfrak{q}$（若 $\mathfrak{q}\in\mathfrak{U}$），否则为 $0$。
（b）证明 $i_!$ 是一个\trExact{}函子。
（c）若 $\mathscr{G}$ 是一个 $\OO_Y$-\trModule，请描述一个包含 $i_!i^{-1}\mathscr{G}\hookrightarrow\mathscr{G}$。（一个有趣的注记，我们用不到：设 $Z$ 为 $\mathfrak{U}$ 的补集，并设 $j\colon Z\to Y$ 为自然包含。那么就存在一个\trShortExactSequence
\[
0 \longrightarrow i_!i^{-1}\mathscr{G} \longrightarrow \mathscr{G} \longrightarrow j_\ast j^{-1}\mathscr{G} \longrightarrow 0.
\]
最好是通过描述这些映射、然后在茎上检验正合性来验证这一点。）
（d）证明 $(i_!,i^{-1})$ 是一个伴随对，从而存在自然双射
\[
\mathrm{Hom}_{\OO_Y}(i_!\mathscr{F},\mathscr{G})\overset{\sim}{\longrightarrow}\mathrm{Hom}_{\OO_\mathfrak{U}}(\mathscr{F},\mathscr{G}|_\mathfrak{U})
\]
对任何 $\OO_\mathfrak{U}$-\trModule{} $\mathscr{F}$ 与 $\OO_Y$-\trModule{} $\mathscr{G}$。（特别地，$\mathscr{G}$ 在 $\mathfrak{U}$ 上的截面可以与 $\mathrm{Hom}_{\OO_Y}(i_!\OO_\mathfrak{U},\mathscr{G})$ 等同起来。）

我们准备好使用\trExtensionByZero{}函子这件工具了。

\textbf{23.4.H.（练习：内射层是松软层）}\quad 设 $(X,\OO_X)$ 是一个\trRingedSpace，而 $\mathscr{Q}$ 是一个内射 $\OO_X$-\trModule。证明 $\mathscr{Q}$ 是\trFlasque{}的。提示：若 $\mathfrak{U}\subset\mathfrak{V}\subset X$，请描述 $\OO_X$-\trModule{}的一个单射 $0\to(i_\mathfrak{U})_!\OO_\mathfrak{U}\to(i_\mathfrak{V})_!\OO_\mathfrak{V}$，其中 $i_\mathfrak{U}\colon\mathfrak{U}\hookrightarrow X$ 与 $i_\mathfrak{V}\colon\mathfrak{V}\hookrightarrow X$ 是显然的开嵌入。对它作用正合的反变函子 $\mathrm{Hom}(\cdot,\mathscr{Q})$。

\textbf{23.4.I.（练习：松软层蕴含 $\Gamma$-\trAcyclic）}\quad 设 $\mathscr{F}$ 是一个\trFlasque{}的 $\OO_X$-\trModule。证明 $\mathscr{F}$ 是 $\Gamma$-\trAcyclic{}的（即对 $i>0$ 有 $\mathrm{H}^i(X,\mathscr{F})=0$，\S23.3.2），做法如下。由于 $\mathit{Mod}_{\OO_X}$ 有足够多内射对象，选取 $\mathscr{F}$ 进入某个内射 $\mathscr{I}$ 的一个包含，并称其余核为 $\mathscr{G}$：
\[
0 \longrightarrow \mathscr{F} \longrightarrow \mathscr{I} \longrightarrow \mathscr{G} \longrightarrow 0.
\]
那么由\mbox{练习~23.4.H}，$\mathscr{I}$ 是\trFlasque{}的；于是由\mbox{练习~23.4.D}(b)，$\mathscr{G}$ 是\trFlasque{}的。取（导出函子）上同调的\trLongExactSequence，证明 $\mathrm{H}^1(X,\mathscr{F})=0$。你的论证对\emph{任何}\trFlasque{}层 $\mathscr{F}$ 都有效，故 $\mathrm{H}^1(X,\mathscr{G})=0$ 也成立。证明 $\mathrm{H}^2(X,\mathscr{F})=0$。把这一点变成一次归纳。

于是，若 $\pi\colon X\to Y$ 是\trRingedSpace{}之间的一个\trMorphism，而 $\mathscr{I}$ 是一个内射 $\OO_X$-\trModule，则 $\mathscr{I}$ 是\trFlasque{}的（\mbox{练习~23.4.H}），于是 $\pi_\ast\mathscr{I}$ 是\trFlasque{}的（\mbox{练习~23.4.E}(b)），于是 $\pi_\ast\mathscr{I}$ 对函子 $\Gamma$ 而言是\trAcyclic{}的（\mbox{练习~23.4.I}），这就完成了\trLeray{}谱序列（\mbox{定理~23.4.5}）的证明。\hfill$\square$

\textbf{23.4.J.（练习）}\quad 把\trLeray{}谱序列（\mbox{定理~23.4.5}）推广到处理
\[
(X,\OO_X) \longrightarrow^{\pi} (Y,\OO_Y) \longrightarrow^{\rho} (Z,\OO_Z)
\]
的高阶前推的复合。

\textbf{23.4.8.}\ \extrahard\ \textbf{\trRingedSpace{}上的 $\OO_X$-\trModule{}范畴未必有足够多射影对象。}\quad 与\mbox{定理~23.4.1} 相对照，\trRingedSpace{}上的 $\OO_X$-\trModule{}的范畴未必有足够多射影对象。例如，设 $X$ 为带扎里斯基拓扑的 $\mathbb{P}^1_k$（事实上，我们对 $X$ 需要知道的非常少——只需要知道它不是一个 \trAlexandrov{}空间），但取 $\OO_X$ 为\trConstantSheaf{} $\mathbb{Z}$。我们将会看到 $\mathit{Mod}_{\OO_X}$——也就是说，$X$ 上阿贝尔群层的范畴——没有足够多射影对象。若 $\mathit{Mod}_{\OO_X}$ 有足够多射影对象，那么就会存在一个从射影层出发的满态射 $\psi\colon P\to\mathbb{Z}$。取定一个闭点 $\mathfrak{q}\in X$。我们将证明茎上的映射 $\psi_\mathfrak{q}\colon P_\mathfrak{q}\to\mathbb{Z}_\mathfrak{q}$ 是零映射，从而与 $\psi$ 的满性矛盾。对 $X$ 的每个开子集 $\mathfrak{U}$，把 $\mathbb{Z}_\mathfrak{U}$ 记作 $\mathfrak{U}$ 到 $X$ 的\trExtensionByZero{}（取值于 $\mathbb{Z}$ 的\trConstantSheaf）（\S23.4.7——$\mathbb{Z}_\mathfrak{U}(\mathfrak{V})=\mathbb{Z}$（若 $\mathfrak{V}\subset\mathfrak{U}$），否则 $\mathbb{Z}_\mathfrak{U}(\mathfrak{V})=0$）。对 $\mathfrak{q}$ 的每个开邻域 $\mathfrak{V}$，设 $W$ 是一个严格更小的开邻域。考虑满态射 $\mathbb{Z}_{X-\mathfrak{q}}\oplus\mathbb{Z}_W\to\mathbb{Z}$。由 $P$ 的射影性，满态射 $\psi$ 提升为 $P\to\mathbb{Z}_{X-\mathfrak{q}}\oplus\mathbb{Z}_W$。映射 $P(\mathfrak{V})\to\mathbb{Z}(\mathfrak{V})$ 穿过 $\mathbb{Z}_{X-\mathfrak{q}}(\mathfrak{V})\oplus\mathbb{Z}_W(\mathfrak{V})=0$ 分解，从而必是零映射。于是茎上的映射 $\psi_\mathfrak{q}\colon P_\mathfrak{q}\to\mathbb{Z}_\mathfrak{q}$ 也是零映射（你看出为什么了吗？），正如所愿。

